%\baselineskip11.8pt plus.2pt minus.2pt


%\normalsize
%\vskip\baselineskip
\advance\baselineskip by.2pt

\niebowspis{Niebo w~czerwcu}

\niebowtyt{Niebo w~czerwcu}

%\vspace*{-4pt}
{
\mathcode`\,="013B
  
Słońce w~czerwcu dociera do szczytu ekliptyki, osiągając 21 dnia miesiąca jej najbardziej na północ wysunięty punkt. W~kolejnych dniach zacznie ono trwającą pół roku wędrówkę na południe. Na początku lata ma miejsce najwcześniejszy wschód i~najpóźniejszy zachód Słońca. Ze względu jednak na to, że Ziemia jest blisko aphelium swojej orbity i~porusza się wolniej niż na początku zimy, to odległość czasowa między tymi zdarzeniami jest dużo mniejsza. Na naszych szerokościach geograficznych Słońce wschodzi najwcześniej 17 czerwca, a~zachodzi najpóźniej 24 czerwca, czyli już tydzień po najwcześniejszym wschodzie. Na czerwiec przypada środek sezonu na obłoki srebrzyste i~łuk okołohoryzontalny,
%\footnote{więcej o~nim na angielskiej stronie: {\tt https://www.atoptics.co.uk/cha2.htm}}
 gdyż Słońce z~jednej strony wędruje najwyżej nad widnokręgiem i~dłużej przebywa na odpowiedniej wysokości, a~z~drugiej strony chowa się pod horyzont najpłycej, i~obłoki srebrzyste można dostrzec na terenie całego kraju.

\vspace*{\parskip}
\begin{dwieszpalty}
Na przełomie wiosny i~lata wieczorem ekliptyka wyraźnie zmniejsza swoje nachylenie do widnokręgu, a~na niebie porannym dopiero zaczyna zwiększać. Dlatego czerwiec nie jest korzystnym miesiącem na obserwacje obiektów przebywających jednocześnie blisko niej i~Słońca. Odbija się to przede wszystkim na widoczności planety Wenus, która 5 dnia miesiąca osiągnie maksymalną elongację zachodnią, oddalając się prawie na $46^{\circ}$ od Słońca. Mimo tego o~świcie zdąży się wtedy wznieść na wysokość zaledwie około $5^{\circ}$ ponad wschodni widnokrąg. Potem, wraz ze zwiększającym się nachyleniem ekliptyki, widoczność planety zacznie się poprawiać. Obserwatorzy nieba na pewno pamiętają, że na początku roku wieczorem, przy bardzo dobrym nachyleniu ekliptyki, planeta wznosiła się zdecydowanie wyżej. Wenus zacznie miesiąc w~Rybach i~przejdzie przez Barana, by do końca czerwca dotrzeć na odległość~$9^{\circ}$ na południowy zachód od Plejad. W~tym czasie jej jasność spadnie do $-4,1^m$, średnica tarczy~-- do~$18''$, faza natomiast urośnie do $63\%$. 22~czerwca $6^{\circ}$~nad planetą pokaże się Księżyc w~fazie $16\%$.

Srebrny Glob rozświetli swoim blaskiem przede wszystkim pierwszą część miesiąca, zaczynając go w~fazie $37\%$ w~Lwie, w~połowie drogi między Regulusem, najjaśniejszą gwiazdą tej konstelacji, a~Marsem, który jest jedyną planetą widoczną w~pierwszej części nocy. Jednak jego warunki obserwacyjne w~czerwcu wyraźnie się pogarszają. Mars przez cały miesiąc pokona ponad $16^{\circ}$, przechodząc w~dniach 16 i~17 czerwca mniej niż $1^{\circ}$ od Regulusa. Na początku planeta około godziny 23 zajmuje pozycję na wysokości prawie $20^{\circ}$ nad zachodnią częścią nieboskłonu, by do końca czerwca o~tej samej porze obniżyć ją do mniej niż $5^{\circ}$. Jasność planety wyniesie około $+1,5^m$, przy średnicy tarczy $5''$. 29~czerwca Księżyc ponownie przejdzie między Marsem a~Regulusem, tym razem jednak planeta znajdzie się po wschodniej stronie gwiazdy, a~faza Srebrnego Globu spadnie do $22\%$.

3 czerwca Księżyc przejdzie przez I~kwadrę, 11~czerwca zaś -- przez pełnię. W~ciągu tych dni odwiedzi obszar od Lwa poprzez Pannę i~Wagę do pogranicza gwiazdozbiorów Skorpiona, Wężownika i~Strzelca. Naturalny satelita Ziemi, przechodząc przez południową część swojej orbity, nadal przebywa jednocześnie głęboko na południe od ekliptyki, co sprawi, że w~okolicach pełni przetnie południk lokalny na wysokości jedynie około $10^{\circ}$. W~nocy z~6 na 7~czerwca Księżyc w~fazie $83\%$ pokaże się około $4^{\circ}$ od Spiki, najjaśniejszej gwiazdy Panny, a~trzy dni później zakryje gwiazdę $\pi$ Scorpii, czyli najbardziej na południe wysuniętą gwiazdę z~charakterystycznego łuku trzech jaśniejszych od $+3^m$ gwiazd na północny zachód od Antaresa, najjaśniejszej gwiazdy Skorpiona. W~Polsce zjawisko potrwa mniej więcej od godziny 22:20 do 23:40, w~zależności od położenia. Im bardziej na północny zachód, tym wcześniej. Niestety wynosząca $98\%$ faza Księżyca znacznie utrudni jego obserwacje.

Po pełni Srebrny Glob podąży ku ostatniej kwadrze -- 18 czerwca i~nowiu -- 25 czerwca. Jednak ze względu na wspomniane niekorzystne nachylenie ekliptyki do porannego widnokręgu zacznie on wschodzić po północy (szczególnie po ostatniej kwadrze), zajmując niskie położenie nad horyzontem. Między pełnią a~ostatnią kwadrą warto wspomnieć o~przejściu w~nocy z~12 na 13~czerwca Księżyca w~fazie~$97\%$ $3^{\circ}$~na~południe od Nunki, jednej z~jaśniejszych gwiazd Strzelca, oraz~spotkaniu z~Saturnem 6 dni później, gdy w~połowie oświetlony Księżyc przejdzie $2^{\circ}$ od planety z~pierścieniami. Około godziny 2 oba ciała niebieskie zdążą się wznieść na wysokość około $10^{\circ}$. Saturn świeci blaskiem $+1^m$, a~jego tarcza ma średnicę $17''$.

Jak co roku, pod koniec miesiąca promieniują meteory z~roju Bootydów, osiągając maksimum aktywności 27~czerwca. Radiant roju znajduje się w~północnej części Wolarza, prawie w~tym samym miejscu co styczniowych Kwadrantydów, i~wieczorem zajmuje pozycję na wysokości ponad $60^{\circ}$ nad zachodnią częścią nieboskłonu. Są to jedne z~wolniejszych meteorów, ich prędkość zderzenia z~naszą atmosferą wynosi zaledwie 18~km/s. W~tym roku nie prognozuje się zwiększonej aktywności tego roju, bliski zaś nowiu Księżyc nie przeszkodzi w~ich obserwacjach.

\medskip
\rightline{\large
  \it  Ariel MAJCHER}
\end{dwieszpalty}

}\eject
\normalsize